\scanpage{211}
\mspage{160}

respectivement, et la base l'est par hypothèse. Il en résulte que la face supérieure du cube est également un carré cartésien.

Ceci vu, \deleted{\unsure{com}} \unsure{comme} $q$ est lisse, il en est de $=$ de $q_{/t}$ dans le carré cartésien $(\unsure{\alpha|a|*})$, donc aussi de $p_{/t}$ par changement de base. D'autre part, la lissité de $q$ implique que $q_{/t}$ est à fibres $W$-asphériques, donc par changement de base il en est de $=$ de $p_{/t}$ (qui est donc une $W$-équivalence universelle, i.e. une $W$-fibration triviale). De plus, \inserted{à cause de la lissité de $p_{/t}$,} pour $v'$ dans $U'_{/t'}$, l'inclusion

\[
(U_{/t})_{v'}\longrightarrow v'\backslash(U^{\deleted{\unsure{\prime}}}_{/t\deleted{\unsure{\prime}}})
\]

est dans $W$, donc (par $W(5)$) définit une bijection sur les $\unsure{\pi_0|\Pi_0}$. Comme $(U_{/t})_{v'}$ est $W$-asphérique par ce qui précède, il s'ensuit \inserted{(par $W(5)$)} que son $\pi_0$ est de cardinal 1, donc il en est de $=$ de $\unsure{v'(U_{/t})|v'\backslash(U_{/t})}$, i.e. celui-ci est $0$-connexe. Or c'est là la condition d) de la prop. \unsure{15|10} qu'il fallait établir.
% archive sheet number: 211
